%!TEX program = pdflatex
\documentclass[aps,%
twocolumn,%
floatfix,%
pra,%
amsfonts,%
longbibliography,%
groupedaddress,%
superscriptaddress]{revtex4-2}%

%%%%%%%%%%%%%%%%%%%%%%%%%%%%%%%%%%%%%%%%%%%%%%%%%%%%%%%%%%%%%%%%%%%%%%%%%
% Packages
%%%%%%%%%%%%%%%%%%%%%%%%%%%%%%%%%%%%%%%%%%%%%%%%%%%%%%%%%%%%%%%%%%%%%%%%%
\usepackage{amsmath}
\usepackage{amssymb}
\usepackage{amsthm}
\usepackage{bm,bbm} % to use the command \mathbbm{1}, \bm{x}
\usepackage{graphicx}
\usepackage{indentfirst}
\usepackage{enumitem}
% \usepackage{algorithmicx}
% \usepackage{algpseudocode}
\usepackage{algorithm}
% \usepackage{algorithmic}
% \usepackage{algcompatible}
\usepackage{algpseudocode}
\usepackage{tikz}
\usepackage{quantikz}
\usetikzlibrary{graphs} 
\usepackage[titletoc,toc,page,title]{appendix}
\definecolor{beamer@blendedblue}{rgb}{0.2,0.2,0.7}      % color used in beamer
\usepackage[bookmarks=false]{hyperref}
\hypersetup{colorlinks=true,%
citecolor=beamer@blendedblue,%
linkcolor=beamer@blendedblue,%
urlcolor=beamer@blendedblue,%
pdfstartview=FitH,%
bookmarksopen=true}

\def\UrlBreaks{\do\/\do-} % url break
\usepackage[T1]{fontenc}
\usepackage{times}
\usepackage{booktabs} % to use the \toprule \midrule command
%%%%%%%%%%%%%%%%%%%%%%%%%%%%%%%%%%%%%%%%%%%%%%%%%%%%%%%%%%%%%%%%%%%%%%%%%
% General Commands
%%%%%%%%%%%%%%%%%%%%%%%%%%%%%%%%%%%%%%%%%%%%%%%%%%%%%%%%%%%%%%%%%%%%%%%%%
\newtheorem{theorem}{Theorem}
\newtheorem{proposition}[theorem]{Proposition}
\newtheorem{lemma}{Lemma}
\newtheorem{corollary}{Corollary}

%%%%%%%%%%%%%%%%%%%%%%%%%%%%%%%%%%%%%%%%%%%%%%%%%%%%%%%%%%%%%%%%%%%%%%%%%
\usepackage{pifont}

%%%%%%%%%%%%%%%%%%%%%%%%%%%%%%%%%%%%%%%%%%%%%%%%%%%%%%%%%%%%%%%%%%%%%%%%%
% Quantum Commands
%%%%%%%%%%%%%%%%%%%%%%%%%%%%%%%%%%%%%%%%%%%%%%%%%%%%%%%%%%%%%%%%%%%%%%%%%
\DeclareMathOperator{\Tr}{Tr}
\DeclareMathOperator{\tr}{Tr}
\DeclareMathOperator{\NOT}{NOT}
\DeclareMathOperator{\diag}{diag}
\DeclareMathOperator{\rank}{rank}
\DeclareMathOperator{\conv}{conv}
\DeclareMathOperator{\sgn}{sgn}
\DeclareMathOperator{\Sp}{Span}
\DeclareMathOperator{\id}{id}
% \newcommand{\bra}[1]{\langle #1\rvert}
% \newcommand{\ket}[1]{\lvert #1\rangle}
% \newcommand{\proj}[1]{\lvert #1\rangle\!\langle #1\rvert}
\newcommand{\ox}{\otimes}
\newcommand{\1}{\mathbbm{1}}
\newcommand{\mean}[1]{\langle #1\rangle}
\newcommand{\gen}[1]{\langle #1\rangle}
\newcommand{\ketbra}[2]{\lvert #1\rangle\!\langle #2\rvert}
% \newcommand{\braket}[2]{\langle #1\vert #2\rangle}
%%%%%%%%%%%%%%%%%%%%%%%%%%%%%%%%%%%%%%%%%%%%%%%%%%%%%%%%%%%%%%%%%%%%%%%%%
%%% Mathematical symbols
%%%%%%%%%%%%%%%%%%%%%%%%%%%%%%%%%%%%%%%%%%%%%%%%%%%%%%%%%%%%%%%%%%%%%%%%%
\newcommand*{\cF}{\mathcal{F}}
\newcommand*{\cH}{\mathcal{H}}
\newcommand*{\cO}{\mathcal{O}}
\newcommand*{\cP}{\mathcal{P}}
\newcommand*{\cS}{\mathcal{S}}
\newcommand*{\bF}{\mathbb{F}}
\newcommand*{\bP}{\mathbb{P}}

% Colors
\definecolor{alizarin}{rgb}{0.82, 0.1, 0.26}
\definecolor{googleblue}{HTML}{4285F4}
\definecolor{googlered}{HTML}{DB4437}
\definecolor{googleyellow}{HTML}{F4B400}
\definecolor{googlegreen}{HTML}{0F9D58}
\definecolor{klevinblue}{HTML}{002FA7}
\definecolor{tiffanyblue}{HTML}{0ABAB5}

\newcommand{\KW}[1]{{\color{googlered}(KW: #1)}}
\newcommand{\SYC}[1]{\textcolor{blue}{(SYC: #1)}}

\begin{document}

\title{Meeting 01: Technical Details}
%%%
\author{Siyuan Chen}
\author{Zhongxia Shang}

%%%


\begin{abstract}
\noindent 
We propose a method to establish criteria for determining whether the effects of noise "obscure" fixed or random circuit based on spectrum of Pauli decomposition, thereby enabling a unified framework for truncation based simulation.
\end{abstract}
\date{\today}
\maketitle

\section{Necessary condition of approximation algorithm}
Given an operator in the form of Pauli decomposition $O = \sum_{i=1}^{N} c_i P_i$ with $N=O(2^n)$ , truncation algorithm with $T =O(\tilde N)$ time complexity calculates an approximation $\tilde O = \sum_{i = 1}^{\tilde N} \tilde c_i \tilde P_i $ with $\tilde N = \text{poly}(n)$ and  try to approximate operator average $\text{tr}(\rho O)$ within precision 
\begin{align}
     \Delta &= \|\text{tr}(\rho O) - \text{tr}(\rho \tilde O)\|^2 \\
    &\leq \|O - \tilde O\|_F^2 := \varepsilon^2,
\end{align}
where we use the random circuit condition in ~\cite{schuster2024polynomial}. For fixed $\tilde N$, the optimal truncation strategy is keeping the largest $\tilde N$ Pauli operators and truncates the others. Supposed that $\{c_i\}_{i = 1}^{N}$ is in the decreasing order with respect to $i$,  this gives the optimal error given fixed $\tilde N$ as $  \varepsilon^2 = \sum_{i = \tilde N}^N c_i^2 $. Noted that such an optimal precision might hard to reach due to the difficulty in selecting $\tilde N$ Pauli operators with the largest $c_i$. Thus we define a function that maps the algorithm complexity to optimal algorithm precision
\begin{align}
    f(\tilde N) = \sum_{i = \tilde N}^N c_i^2.
\end{align}
\begin{figure}[htbp!]
    \centering
    \includegraphics[width=1\linewidth]{pic/necessary_condition.png}
    \caption{Noiseless circuit decay slowly compare to noisy circuit. Effective truncation method exist only if the decay of $c_i^2$ is faster than $i^{-1}$}
    \label{fig:necessary_condition}
\end{figure}
If the optimal precision transformed as $\varepsilon \to \varepsilon/A$, the complexity of the algorithm should at least scales as $p(A)$, which is a polynomial of $A$ with rank $r$. In other words,
\begin{align}
    f[\tilde N p(A)] &= \frac{1}{A^2} f(\tilde N)
\end{align}
Take the derivative of $A$ gives the identity
\begin{align}
    c^2_{\tilde N p(A)} = \frac{2 f(\tilde N)}{A^3  \tilde N p'(A)}.
\end{align}
When $A$ is large enough, we conclude that $p(A)= \Theta(A^r), p'(A)= \Theta(A^{r-1})$, This indicates that the constant $c_i$ must decrease faster than $i^{-1}$
\begin{align}
    c_{i}^2 = \Theta(i^{-1 -2/r}) \leq \Theta(i^{-1}).
\end{align}

As a concrete example, we noted that in reference~\cite{schuster2024polynomial}, it is proved that the optimal error satisfied $\varepsilon^2 \leq e^{-\gamma(l+1)} \cdot \|O\|_F^2$. In the main time, the time complexity satisfied $T = O(n^{l/d})$. Thus if we increase the precision by $l \to l+ l_0$, we have
\begin{align}
    \varepsilon^2 \to e^{-\gamma l_0}\varepsilon, \quad T \to T\cdot n^{l_0/d}
\end{align}
\section{Open question and Research directions}


%%%%%%%%%%%%%%%%%%%%%%%%%%%%%%%%%%%%%%%%%%%%%%%%%%%%%%%%%%%%%%%%%%%%%%%%%
%%%%%%%%%%%%%%%%%%%%%%%%%%%%%%%%%%%%%%%%%%%%%%%%%%%%%%%%%%%%%%%%%%%%%%%%%
\end{document}